\documentclass[UTF8,a4paper,12pt]{ctexart}

\usepackage{geometry}
\usepackage{amsmath}
\usepackage{amssymb}
\usepackage{booktabs}
\usepackage{array}
\usepackage{longtable}
\usepackage{enumitem}
\usepackage{hyperref}
\usepackage{xcolor}

\geometry{left=2.5cm,right=2.5cm,top=2.4cm,bottom=2.4cm}
\hypersetup{
  colorlinks=true,
  linkcolor=blue,
  urlcolor=blue
}
\setlist[itemize]{leftmargin=2em}
\setlist[enumerate]{leftmargin=2em}

\title{共享单车再平衡调度优化模型说明}
\author{深圳共享单车再平衡项目}
\date{\today}

\begin{document}
\maketitle

\section{建模目标与基本思路}

本项目希望在给定若干调度车辆、区域单车富余量、区域单车短缺量和区域间距离的条件下，设计共享单车再平衡调度方案。与只最小化调度距离的模型不同，本文模型将不同调度车同时行动、车辆容量差异、装卸单位单车所需时间以及未满足短缺惩罚同时纳入目标函数。

模型的核心目标采用改进后的复合目标函数：

\[
\min \quad
\alpha T+\gamma\sum_{k\in K}T_k+\sum_{j\in D}\beta_j u_j
\]

其中，\(T\) 表示所有调度车辆并行行动后的总完成时间，即最晚完成任务车辆的完成时间；\(T_k\) 表示车辆 \(k\) 的服务时间；\(u_j\) 表示短缺区域 \(j\) 未被满足的单车数量；\(\alpha\)、\(\gamma\) 和 \(\beta_j\) 分别表示最大完成时间、总服务时间和区域短缺惩罚的权重。

该目标函数体现三个优化方向：

\begin{itemize}
  \item 尽量缩短整个调度任务完成所需时间；
  \item 控制所有车辆的总运营时间，避免为了均衡车辆任务而产生过多绕行；
  \item 尽量减少未被满足的短缺车辆数量。
\end{itemize}

由于不同调度车可以同时出发并并行执行任务，因此不宜只最小化所有车辆服务时间之和。若只最小化总服务时间，模型可能让一辆车承担很长路线，而其他车辆闲置，这与“多车并行、尽快完成调度”的目标不一致。因此，本文使用最大完成时间 \(T=\max_k T_k\) 表示整体调度完成时间。

\section{集合、参数与变量}

\subsection{集合}

\begin{longtable}{>{\centering\arraybackslash}p{3cm}p{11cm}}
\toprule
符号 & 含义 \\
\midrule
\endhead
\(S\) & 富余区域集合，区域内可调出的单车数量大于 0 \\
\(D\) & 短缺区域集合，区域内需要补入的单车数量大于 0 \\
\(N=S\cup D\) & 所有需要考虑的服务区域集合 \\
\(E\) & 稀疏有向邻接边集合，只有 \((i,j)\in E\) 时车辆才允许从 \(i\) 行驶到 \(j\) \\
\(K\) & 调度车辆集合 \\
\(0\) & 调度中心、车辆初始点或统一起终点 \\
\bottomrule
\end{longtable}

在项目数据中，区域可以采用 1km 网格或 500m 网格。若使用 1km 网格，则 \(S\) 和 \(D\) 可由 `scenario\_grid\_demand\_1km.parquet` 中的 `surplus` 和 `shortage` 字段得到。

为降低模型规模，本文不再采用全联通图，而是根据区域距离构造稀疏邻接图。例如：

\[
E=\{(i,j):d_{ij}\le R,\ i\ne j\}
\]

或者令每个区域仅连接其最近的 \(k_e\) 个邻居。若某些区域因此无法从调度中心到达，则为其补充最近邻边，直到稀疏图满足基本连通性要求。

\subsection{参数}

\begin{longtable}{>{\centering\arraybackslash}p{3cm}p{11cm}}
\toprule
符号 & 含义 \\
\midrule
\endhead
\(d_{ij}\) & 区域 \(i\) 到区域 \(j\) 的距离，可由网格中心点距离近似得到 \\
\(R\) & 稀疏邻接图连接半径，只有距离不超过 \(R\) 的区域对才默认建边 \\
\(k_e\) & 每个区域至少保留的最近邻数量，用于避免局部边过少导致图不连通 \\
\(v_k\) & 调度车 \(k\) 的平均速度；若速度相同，可令 \(v_k=v\) \\
\(Q_k\) & 调度车 \(k\) 的最大装载容量，不同车辆可以不同 \\
\(a\) & 装载一辆单车所需时间 \\
\(b\) & 卸载一辆单车所需时间 \\
\(s_i\) & 富余区域 \(i\) 可被调走的单车数量，\(i\in S\) \\
\(r_j\) & 短缺区域 \(j\) 需要补入的单车数量，\(j\in D\) \\
\(\alpha\) & 最大完成时间成本权重 \\
\(\gamma\) & 车辆总服务时间成本权重 \\
\(\beta_j\) & 区域 \(j\) 的单位未满足短缺惩罚权重，可反映区域重要性 \\
\(M\) & 足够大的常数，用于 Big-M 约束 \\
\bottomrule
\end{longtable}

\subsection{决策变量}

\begin{longtable}{>{\centering\arraybackslash}p{3cm}p{11cm}}
\toprule
变量 & 含义 \\
\midrule
\endhead
\(x_{ijk}\in\{0,1\}\) & 若车辆 \(k\) 沿稀疏边 \((i,j)\in E\) 行驶，则取 1，否则取 0 \\
\(y_{ik}\in\{0,1\}\) & 若车辆 \(k\) 访问或经过区域 \(i\)，则取 1，否则取 0 \\
\(h_{ik}\in\{0,1\}\) & 若车辆 \(k\) 在区域 \(i\) 发生装车或卸车服务，则取 1，否则取 0 \\
\(p_{ik}\ge 0\) & 车辆 \(k\) 在富余区域 \(i\) 装载的单车数量 \\
\(q_{jk}\ge 0\) & 车辆 \(k\) 在短缺区域 \(j\) 卸载的单车数量 \\
\(L_{ik}\ge 0\) & 车辆 \(k\) 离开区域 \(i\) 时车上的单车数量 \\
\(u_j\ge 0\) & 短缺区域 \(j\) 未被满足的单车数量 \\
\(T_k\ge 0\) & 车辆 \(k\) 的服务完成时间 \\
\(T\ge 0\) & 所有车辆并行调度后的总完成时间，即最大车辆完成时间 \\
\(o_{ik}\ge 0\) & MTZ 子回路消除变量，表示车辆 \(k\) 访问区域 \(i\) 的顺序 \\
\bottomrule
\end{longtable}

\section{数学模型}

\subsection{目标函数}

\[
\min \quad
\alpha T+\gamma\sum_{k\in K}T_k+\sum_{j\in D}\beta_j u_j
\]

该目标同时考虑整体调度完成时间、车辆总运营时间和区域短缺惩罚。其中，\(\alpha T\) 体现多车并行条件下的最终完成时间；\(\gamma\sum_k T_k\) 控制总体运营成本；\(\sum_j\beta_j u_j\) 则惩罚未满足短缺。若核心区域的短缺影响更大，可为这些区域设置更高的 \(\beta_j\)。

\subsection{车辆完成时间约束}

每辆调度车的服务时间由三部分构成：行驶时间、装车时间和卸车时间。

\[
T_k =
\sum_{(i,j)\in E}
\frac{d_{ij}}{v_k}x_{ijk}
+a\sum_{i\in S}p_{ik}
+b\sum_{j\in D}q_{jk}
+c\sum_{i\in N}h_{ik},
\quad \forall k\in K
\]

其中 \(c\) 表示一次真实停靠服务的固定时间。车辆可以经过中间网格，但只有当 \(h_{ik}=1\) 即发生装卸服务时，才计入固定停车、找车和操作时间。

由于调度车同时行动，整体完成时间应不小于任意车辆的完成时间：

\[
T\ge T_k,\quad \forall k\in K
\]

\subsection{富余区域约束}

每个富余区域被调走的单车数量不能超过该区域的富余量：

\[
\sum_{k\in K}p_{ik}\le s_i,\quad \forall i\in S
\]

车辆只有在某个富余区域发生服务，才能在该区域装车：

\[
p_{ik}\le s_i h_{ik},\quad \forall i\in S,\ k\in K
\]

\subsection{短缺区域约束}

短缺区域可以不被完全满足，但未满足部分需要进入目标函数的惩罚项：

\[
\sum_{k\in K}q_{jk}+u_j=r_j,\quad \forall j\in D
\]

车辆只有在某个短缺区域发生服务，才能在该区域卸车：

\[
q_{jk}\le r_j h_{jk},\quad \forall j\in D,\ k\in K
\]

服务节点必须被车辆路径访问：

\[
h_{ik}\le y_{ik},\quad \forall i\in N,\ k\in K
\]

\subsection{车辆容量与载重变化约束}

不同车辆具有不同装载容量：

\[
0\le L_{ik}\le Q_k,\quad \forall i\in N\cup\{0\},\ k\in K
\]

车辆离开调度中心时可设为空车：

\[
L_{0k}=0,\quad \forall k\in K
\]

如果车辆 \(k\) 从区域 \(i\) 行驶到区域 \(j\)，则其离开 \(j\) 时的载重应等于离开 \(i\) 时的载重加上在 \(j\) 的装车量、减去在 \(j\) 的卸车量。采用 Big-M 形式可写为：

\[
L_{jk}\ge L_{ik}+p_{jk}-q_{jk}-M(1-x_{ijk})
\]

\[
L_{jk}\le L_{ik}+p_{jk}-q_{jk}+M(1-x_{ijk})
\]

其中 \((i,j)\in E\)、\(k\in K\)。对于非富余点，可令 \(p_{jk}=0\)；对于非短缺点，可令 \(q_{jk}=0\)。

\subsection{路径流平衡约束}

若车辆访问某个区域，则必须有一条进入该区域的路径和一条离开该区域的路径：

\[
\sum_{j:(i,j)\in E}x_{ijk}=y_{ik},
\quad \forall i\in N,\ k\in K
\]

\[
\sum_{j:(j,i)\in E}x_{jik}=y_{ik},
\quad \forall i\in N,\ k\in K
\]

每辆车最多从调度中心出发一次，并最多返回调度中心一次：

\[
\sum_{j:(0,j)\in E}x_{0jk}\le 1,\quad \forall k\in K
\]

\[
\sum_{i:(i,0)\in E}x_{i0k}\le 1,\quad \forall k\in K
\]

若要求所有车辆必须参与调度，可将上述不等式改为等式。但在实际场景中，允许车辆不出车通常更合理，因为当短缺量较少或惩罚较低时，强制所有车辆出车会产生不必要的时间成本。

\subsection{子回路消除：MTZ 约束}

仅有流平衡约束时，模型可能生成不经过调度中心的小回路。例如，某辆车的路径中可能同时出现 \(i\to j\to i\) 这样的独立环路，但该环路并未与调度中心相连。为避免这种情况，本文采用 MTZ 约束消除子回路。

为每辆车 \(k\) 和每个服务区域 \(i\in N\) 设置访问顺序变量 \(o_{ik}\)。若车辆 \(k\) 不访问区域 \(i\)，则顺序变量为 0；若访问，则顺序变量位于 \(1\) 到 \(|N|\) 之间：

\[
0\le o_{ik}\le |N|y_{ik},\quad \forall i\in N,\ k\in K
\]

对任意两个不同服务区域 \(i,j\in N\)，加入：

\[
o_{ik}-o_{jk}+|N|x_{ijk}\le |N|-1,
\quad \forall (i,j)\in E,\ i,j\in N,\ k\in K
\]

当 \(x_{ijk}=1\) 时，上式等价于要求车辆访问 \(j\) 的顺序必须晚于访问 \(i\) 的顺序，从而阻止服务区域之间形成与调度中心无关的闭合子环。该方法形式简单，适合课程项目中的混合整数规划建模。

\section{模型结果可能与目标不符的风险}

虽然上述模型比简单的两阶段运输模型更贴近真实调度，但仍存在一些结果与实际目标不完全一致的风险。下面逐项分析。

\subsection{风险一：权重设置不当导致偏向极端解}

目标函数为：

\[
\min \alpha T+\gamma\sum_{k\in K}T_k+\sum_{j\in D}\beta_j u_j
\]

如果 \(\beta_j\) 过大，模型会为了减少少量短缺而安排车辆行驶很远，导致调度时间过长；如果 \(\beta_j\) 过小，模型可能宁愿不补车，从而产生较大的未满足短缺。若 \(\alpha\) 远大于 \(\gamma\)，模型会过度关注最晚完成时间；若 \(\gamma\) 远大于 \(\alpha\)，模型会更接近最小化总运营时间，可能牺牲并行完成效率。

\textbf{改进建议：}

\begin{itemize}
  \item 对 \(\alpha\)、\(\gamma\) 和 \(\beta_j\) 做敏感性分析，比较不同权重下的最大完成时间、总服务时间和短缺量。
  \item 若业务上要求尽量先满足短缺，可采用分层目标：先最小化未满足短缺，再在短缺最小的基础上最小化完成时间。
  \item 可设置最低服务水平约束，例如 \(\sum_j u_j \le \eta \sum_j r_j\)，限制总未满足比例。
\end{itemize}

\subsection{风险二：只惩罚总短缺量，可能忽视区域公平性}

若所有区域使用相同惩罚权重，模型可能集中服务容易到达的大短缺点，而长期忽略偏远小短缺区域。总短缺量看起来较小，但服务分布可能不公平。

\textbf{改进建议：}

\begin{itemize}
  \item 保留区域化惩罚权重 \(\beta_j\)，将核心商圈、地铁站等重点区域赋予更高权重。
  \item 加入区域最大未满足比例约束：\(u_j\le \rho r_j\)。
  \item 在报告中区分“总短缺最小”和“服务公平性”，避免将两者混为一谈。
\end{itemize}

\subsection{风险三：最大完成时间与总运营时间存在权衡}

目标函数已经加入 \(\gamma\sum_k T_k\) 来控制总运营时间，但如果 \(\alpha\) 远高于 \(\gamma\)，模型仍可能为了缩短最晚完成时间而让部分车辆绕路；如果 \(\gamma\) 远高于 \(\alpha\)，模型则可能减少总工作量，却让某一辆车承担较长任务。

\textbf{改进建议：}

\begin{itemize}
  \item 通过敏感性分析选择 \(\alpha\) 与 \(\gamma\) 的比例。
  \item 若公司更关心总人力和油耗成本，则提高 \(\gamma\)。
  \item 若更关心高峰前快速完成，则提高 \(\alpha\)。
\end{itemize}

\subsection{风险四：距离近似可能低估实际服务时间}

若使用网格中心球面距离 \(d_{ij}\) 近似道路行驶距离，则模型可能低估跨河、主干路绕行或交通拥堵造成的实际时间差异。这样得到的最优路径在地图上看似合理，但实际执行可能耗时更长。

\textbf{改进建议：}

\begin{itemize}
  \item 若能获取路网数据，应使用道路网络最短路时间替代直线距离。
  \item 若无法获取路网数据，可将距离乘以经验修正系数。
  \item 可对不同行政区、主干道隔离区域设置额外时间惩罚。
\end{itemize}

\subsection{风险五：装卸时间采用线性假设可能过于简单}

模型假设装一辆车花费 \(a\)，卸一辆车花费 \(b\)，总装卸时间与车辆数量线性相关。但实际中，停车、寻找车辆、扫码、搬运、装车整理等存在固定时间成本。

\textbf{改进建议：}

\begin{itemize}
  \item 采用服务停靠变量 \(h_{ik}\)，只在发生装卸时计入固定服务时间。
  \item 将车辆装卸时间表示为“固定服务时间 + 单位车辆装卸时间”。
  \item 对大规模装卸和小规模装卸设置不同效率参数。
\end{itemize}

\subsection{风险六：允许部分短缺未满足可能产生不可接受方案}

引入 \(u_j\) 可以避免模型因车辆不足或时间成本过高而不可行，但如果没有服务水平下限，模型可能选择大量不服务，只支付惩罚。

\textbf{改进建议：}

\begin{itemize}
  \item 设置总满足率约束：\(\sum_j (r_j-u_j)\ge \theta \sum_j r_j\)。
  \item 对重点区域设置强制满足约束：\(u_j=0\) 或 \(u_j\le \rho_j r_j\)。
  \item 把短缺惩罚权重 \(\beta\) 设置为高于合理调度成本的水平。
\end{itemize}

\subsection{风险七：车辆初始位置和是否返场影响结果}

如果假设所有车辆都从同一个调度中心出发并返回，模型较简单，但可能与真实情况不符。实际车辆可能分布在不同仓库或上一轮任务结束点。

\textbf{改进建议：}

\begin{itemize}
  \item 为每辆调度车设置独立起点 \(o_k\) 和终点 \(h_k\)。
  \item 若车辆不需要返场，则删除返回调度中心约束。
  \item 若车辆需要在规定时间后回到仓库，应保留终点约束并计入返回时间。
\end{itemize}

\subsection{风险八：模型规模较大，求解困难}

该模型属于带容量约束的车辆路径问题，且包含装卸数量变量和短缺惩罚，整体上是混合整数规划问题。若网格数量较多、车辆数量较多，精确求解可能非常慢。

\textbf{改进建议：}

\begin{itemize}
  \item 先筛选短缺和富余量较大的区域，减少节点数量。
  \item 使用 1km 网格作为主模型，500m 网格作为敏感性分析。
  \item 采用两阶段启发式：先聚类或匹配，再做局部路径优化。
  \item 设置求解时间上限，使用可行解与下界差距作为结果质量说明。
\end{itemize}

\section{服务水平约束与参数建议}

为了使结果更符合实际服务要求，可以在复合目标函数之外加入服务水平约束。由于实际场景中总富余量可能小于总短缺量，服务水平不应直接以总短缺量为基准，而应以可满足上限为基准：

\[
\sum_{j\in D}(r_j-u_j)\ge \theta \min\left(\sum_{j\in D}r_j,\sum_{i\in S}s_i\right)
\]

等价地：

\[
\sum_{j\in D}u_j
\le
\sum_{j\in D}r_j
-\theta \min\left(\sum_{j\in D}r_j,\sum_{i\in S}s_i\right)
\]

其中 \(\theta\) 为最低满足率。例如，若要求至少满足可满足上限的 90\%，则取 \(\theta=0.9\)。这种写法可以避免在“总富余量小于总短缺量”时把模型错误地变成不可行。

\section{结论}

本文模型将共享单车再平衡问题建模为异质车辆、带容量约束、考虑装卸时间和短缺惩罚的车辆路径优化问题，并直接采用复合目标函数：

\[
\min \quad
\alpha T+\gamma\sum_{k\in K}T_k+\sum_{j\in D}\beta_j u_j
\]

相比先求区域间调度量 \(x_{ij}\)、再规划车辆路径的两阶段方法，该模型能更直接地反映车辆容量、实际访问顺序、多车并行、总运营成本和未满足短缺之间的权衡。

同时，本文明确采用 MTZ 约束消除子回路，使每辆车的访问节点形成从调度中心出发并与调度中心相连的有效路径。不过，模型结果仍可能受到权重设置、距离近似、服务公平性、车辆初始位置、装卸时间假设和模型规模的影响。因此，在课程项目中建议采用以下表述：

\begin{quote}
本模型以最晚车辆完成时间、车辆总服务时间和区域化未满足短缺惩罚为联合优化目标，并通过车辆容量、路径流平衡、装卸数量、服务时间、短缺满足和 MTZ 子回路消除约束描述共享单车再平衡过程。为避免模型结果偏离实际运营目标，进一步引入最低服务水平约束，并通过敏感性分析评估参数选择对结果的影响。
\end{quote}

这样既能体现数学模型的完整性，也能说明模型假设与实际应用之间的差距。

\end{document}
